\section{Stage 0: the GPU walk}
\label{sec:gpuwalk}

\subsection{Invocation model}

The walk dispatch has one invocation per compacted scan cell:
\begin{equation}
  \text{groups} = \left\lceil \frac{|\text{cells}|}{64} \right\rceil,
\end{equation}
with a local size of 64.  The shader starts from
\code{gl_GlobalInvocationID.x}; if it is beyond the cell count, it
returns.  All per-invocation state is private:
\begin{itemize}
\item break list (64 floats),
\item local break list (32 floats),
\item task stack (256 intervals),
\item QEF position memo (8 entries),
\item overflow flag and diagnostic counters.
\end{itemize}
No workgroup shared memory is needed.  This makes the pass portable across
devices with small shared-memory budgets and avoids workgroup
synchronization.

\subsection{Node accessors}

The shader wraps the eight storage buffers in typed accessors:
\begin{lstlisting}[style=glsl,caption={Node accessors (abridged).},label={lst:accessors}]
float nodeSdf(uint id, int i) { return (i < 4) ? sdf0[id][i] : sdf1[id][i-4]; }
int nodeType(uint id)   { return meta[id].y; }
int nodeBrush(uint id)  { return meta[id].z; }
int nodeSlot(uint id)   { return int(meta[id].w >> 1) - 1; }
vec3 nodeMin(uint id)   { return cubeMinSize[id].xyz; }
float nodeSize(uint id) { return cubeMinSize[id].w; }
int nodeChild(uint id, int i) { ... }
\end{lstlisting}
The table is a flat array; there are no pointers.  The accessors are
inlined by the compiler.

\subsection{\code{findCellAt}}

The descent mirrors the CPU walker:
\begin{lstlisting}[style=glsl,caption={The root-to-cell descent.},label={lst:findcell}]
Cell findCellAt(vec3 pos) {
    if (!worldContains(pos)) return INVALID;
    uint idx = pc.fromId >= 0 ? pc.fromId : 0;
    if (pos outside anchor cube) idx = 0;
    while (!nodeIsLeaf(idx) && nodeLod(idx) > pc.targetLod) {
        ivec3 half = nodeMin(idx) + 0.5*nodeSize(idx);
        int child = (pos.x >= half.x ? 4 : 0)
                  + (pos.y >= half.y ? 2 : 0)
                  + (pos.z >= half.z ? 1 : 0);
        int c = nodeChild(idx, child);
        if (c < 0) return INVALID;
        idx = c;
    }
    return Cell(idx, nodeMin(idx), nodeSize(idx));
}
\end{lstlisting}
The stop condition is ``leaf or stored LoD at most target''.  The stored
LoD is +1 shifted, exactly as on the CPU.  Starting at \code{fromId} is
valid only for positions strictly inside the anchor cube; a point on the
upper face must follow the sibling branch, so the test uses a strict upper
bound.

\subsection{Edge span construction}

For each of the twelve cube edges, the shader loads the two corner signs.
If they are equal, the edge is inactive.  Otherwise it builds an
\code{EdgeSpan}.  The varying axis is the component with the largest
absolute corner difference (exact for an axis-aligned cube).  The two other
axes are the fixed $u,v$ axes.  The span's epsilon is
\begin{equation}
  \epsilon = \max(\text{worldSize}, L)\cdot10^{-6},
\end{equation}
and the boundary sample offset is
\begin{equation}
  \delta = \max(L\cdot10^{-3},10^{-6}).
\end{equation}
The epsilon is used for duplicate rejection and sign-change tests; the
offset is used when a sample sits exactly on a cell face.

\subsection{The ulp nudge and the zero guard}

Sampling a point exactly on a cell face is ambiguous: the two cells both
claim it.  The shader nudges the fixed coordinates by one float in the
direction of the quadrant.  At zero, \code{nextafter(0,+inf)} is a
subnormal; FTZ GPUs flush it to zero, and the negative quadrant then
resolves into the positive cell.  The shader detects a subnormal nudge and
uses a small normal offset instead.  This is a small detail with a large
visual consequence: without it, the surface disappears along the world
$x=0$, $y=0$, and $z=0$ planes.

\subsection{The scan loop}

For each active edge the walk:
\begin{enumerate}
\item initializes the global break list with the span endpoints;
\item calls \code{collectBreaks};
\item sorts and uniques the global list;
\item iterates consecutive breaks and calls \code{emitSegment} for each
  interval wider than $2\epsilon$.
\end{enumerate}
This is exactly the CPU loop.  The only difference is that the CPU stores
the breaks in a \code{std::vector} and the GPU uses fixed arrays.

\subsection{Diagnostics}

Two counters are accumulated only when \code{VULKAN_SETUP_DIAG} is set:
\code{findCellCalls} and \code{reuseHits}.  They are summed locally and
published with one atomic pair per scan cell, so the diagnostic path does
not add a per-call atomic.  The normal path has a predictable branch and no
counter traffic.

\begin{figure}[htbp]
\centering
\begin{tikzpicture}[node distance=5mm and 8mm,font=\small]
  \node[box] (cell) {one scan cell};
  \node[box,below=of cell] (edges) {12 edges};
  \node[box,below=of edges] (active) {sign change?};
  \node[box,right=of active] (span) {build \code{EdgeSpan}};
  \node[box,below=of span] (breaks) {\code{collectBreaks}};
  \node[box,below=of breaks] (segments) {emit each sub-segment};
  \node[box,below=of segments] (poly) {owner + ring + winding};
  \node[box,below=of poly] (tri) {emit triangle keys};
  \draw[arrow] (cell)--(edges); \draw[arrow] (edges)--(active);
  \draw[arrow] (active)--node[above,note]{yes}(span);
  \draw[arrow] (span)--(breaks); \draw[arrow] (breaks)--(segments);
  \draw[arrow] (segments)--(poly); \draw[arrow] (poly)--(tri);
  \draw[arrow] (active.west) -- ++(-1.0,0) |- node[pos=0.75,left,note]{no, next edge} (edges.west);
\end{tikzpicture}
\caption{One scan-cell invocation.  Inactive edges are skipped; active
edges are stitched and emitted.  The entire control flow is private to the
invocation.}
\label{fig:scanflow}
\end{figure}
